\documentclass[11pt,a4paper]{article}
\usepackage[T1]{fontenc}
\usepackage[utf8]{inputenc}
\usepackage{lmodern}
\DeclareUnicodeCharacter{1E93}{\d{z}}
\usepackage{amsmath,amssymb}
\usepackage[margin=25mm]{geometry}
\usepackage[hidelinks]{hyperref}
\hypersetup{pdftitle={A thermal action scale from Newton's central-force limit},pdfauthor={navstokgap research working notes}}
\newcommand{\tightlist}{\setlength{\itemsep}{0pt}\setlength{\parskip}{0pt}}
\setlength{\emergencystretch}{3em}
\setcounter{secnumdepth}{0}
\begin{document}
\section{A thermal action scale from Newton's central-force
limit}\label{a-thermal-action-scale-from-newtons-central-force-limit}

\textbf{6.1 Sol Max, 2026-11-19.} A central force proportional to radius
admits a test of zero physical action in Newton's own time. Fix a body
in a plane with positive angular momentum. Prepare its radial degree of
freedom in thermal equilibrium and allow changes of the strength of this
one force. If a positive ensemble of Newtonian trajectories obeys the
equilibrium work identity, doubling the orbital frequency costs exactly
\(k_B\log 2\) in isothermal entropy loss at \emph{every} temperature.
Nernst's heat theorem requires this loss to tend to zero. For the
inverse-square force, two changes of strength separated by a short
interval force the equilibrium radius law by the acceleration term in
Lemma X. With both circular orbits inside fixed reflecting walls, the
entropy loss tends to the logarithm of their radial frequency ratio in
units of \(k_B\). Both exclusions include singular ensembles and survive
arbitrary time subdivision.

The two-temperature work observable below defines a positive thermal
action from the body's equation of state, given the usual
high-temperature correspondence and harmonic virial relation. In the
regular planar Schrödinger family it determines the reduced action \(a\)
uniquely, including a finite-observation lower bound. The independent
physical premise concerns thermal equilibria of this rotating body. The
value across different bodies, and an identification with quantum phase
in other theories, remain open.

\subsection{From the finite force to a thermodynamic
question}\label{from-the-finite-force-to-a-thermodynamic-question}

In 1687 Newton specified a regular force in Lemma X. In 1713 he
specified a finite force, allowed it to vary, and added Corollary 4
relating it to initial displacement divided by time squared. Proposition
VI applies the construction to orbits. Proposition X finds the force
toward the centre of an ellipse proportional to distance, with equal
periods for all its ellipses at a common strength. These are edition
passages. The thermal condition in this note is modern.
\href{newton-insertion-action.md\#lemma-x-and-the-measurement-of-initial-force}{Lemma
X in 1687 and 1713} and
\href{https://www.newtonproject.ox.ac.uk/view/texts/diplomatic/NATP00082}{Newton,
1713, Book I, Proposition X and its corollaries}.

Write the acceleration as \(-\omega^2 {\bf q}\). At a point
\({\bf q}_0\ne0\) the departure from its inertial tangent, divided by
\(t^2/2\), tends to \(-\omega^2{\bf q}_0\). This fixes the frequency in
mechanical time. The independently fixed sector is
\(J=q_xp_y-q_yp_x>0\), with mass \(m>0\). A radial thermal contact can
exchange energy while preserving \(J\). Write \(\theta=k_BT>0\) for
temperature in energy units. At each fixed strength the initial law is
stationary. During a prescribed strength change the body is isolated.
For a Newtonian path the full source work on it is the change in

\[
 H_\omega=K+V=\frac{|{\bf p}|^2}{2m}
                 +\frac{m\omega^2|{\bf q}|^2}{2}.
\]

Two equilibrium conditions have independent content.

\begin{enumerate}
\def\labelenumi{\arabic{enumi}.}
\tightlist
\item
  For an initial equilibrium and each smooth positive strength protocol
  from \(\omega\) to \(\omega_f\), inclusive work satisfies
  \(\langle e^{-W/\theta}\rangle=\exp[-(F(\theta,\omega_f)-
  F(\theta,\omega))/\theta]\), with finite Helmholtz free energies. The
  right side is common to all routes between the two strengths. This is
  Jarzynski's identity for an initially thermal, subsequently isolated
  drive. It is used here as a constraint on \emph{any} candidate law.
  Its usual forward derivation starts with a Gibbs state. The
  corresponding quantum work is the difference of the two endpoint
  energy readings.
  \href{https://arxiv.org/pdf/cond-mat/9610209}{Jarzynski, equations (2)
  and (8)} and \href{https://arxiv.org/pdf/0711.3914}{Deffner and Lutz,
  equations (1)--(3)} (passages).
\item
  At fixed \(J\) assume differentiable equilibrium thermodynamics,
  \(S=-k_B\partial_\theta F\), and the Nernst heat theorem for the two
  strengths. In symbols \(S(\theta,\omega)-S(\theta,2\omega)\to0\) as
  \(\theta\downarrow0\). It concerns entropy differences within the same
  equilibrium sector. Its justification as a thermodynamic postulate is
  independent of an action constant. Applying it to this constrained
  radial system is a physical hypothesis. Independent copies give the
  same test per copy by entropy additivity.
  \href{https://www.nobelprize.org/uploads/2017/03/nernst-lecture2.pdf}{Nernst,
  Nobel lecture, pp.~358--359} and
  \href{https://arxiv.org/pdf/1412.3828}{Masanes and Oppenheim, opening
  discussion of the heat theorem and its degeneracy condition}
  (passages).
\end{enumerate}

\subsection{What a quarter-period change forces on every Newtonian
law}\label{what-a-quarter-period-change-forces-on-every-newtonian-law}

\textbf{Theorem.} At fixed \(\theta,\omega,J>0\), condition 1 forces the
Newtonian energy to have law \(E_0+X\), with \(E_0=\omega J\) and \(X\)
exponential of mean \(\theta\). For every \(c>1\) it forces

\[
 F(\theta,c\omega)-F(\theta,\omega)=(c-1)E_0+\theta\log c.       \tag{1}
\]

No density relative to phase volume is required. A hidden label may
affect the initial law arbitrarily if the projected motion and work
still have the specified mechanical values.

\textbf{Proof.} On a stationary orbit put \(E=K+V\) and \(D=V-K\). The
pair \((D,\omega{\bf q}\cdot{\bf p})\) rotates at \(2\omega\) and has
squared radius \(E^2-E_0^2\). This follows from the harmonic equations
and the planar Gram determinant \(J^2\). Stationarity gives
\(\mathbb E[D\mid E]=0\) and \(\mathbb E[D^2\mid E]=(E^2-E_0^2)/2\),
including the circle. The orbit's orientation need never be specified.

Between the initial and final strengths, use frequency
\(\omega_p=\sqrt c\,\omega e^{s/2}\) for its quarter-period
\(\pi/(2\omega_p)\). Its exact map is
\(({\bf q},{\bf p})\mapsto({\bf p}/(m\omega_p),-m\omega_p{\bf q})\). The
final energy at frequency \(c\omega\) is

\[
 E_f=c(e^s V+e^{-s}K)=c(E\cosh s+D\sinh s).                    \tag{2}
\]

These switches are legitimate limits of the smooth finite forces in
condition 1. For \(|s|<\log c\) one has \(E_f\ge ce^{-|s|}E>E\).
Smoothing each switch over a shrinking time changes the linear
fundamental matrix continuously, so the same strict quadratic-form
bound, with a smaller margin, holds for sufficiently narrow ramps.
Dominated convergence applies to \(e^{-(E_f-E)/\theta}\), bounded by
one. Inserting times into a fixed smooth motion composes its exact maps.
Convergent Newton polygons give the same fundamental matrix at fixed
duration. \href{balaban-ode-lesson.md}{Finite-force polygon limit}.

Condition 1 makes this work average independent of \(s\). Differentiate
twice at \(s=0\). On any smaller neighborhood of zero the integrands and
both derivatives are bounded by a polynomial in \(E\) times a decreasing
exponential, even for an initial law with heavy tails. With
\(t=(c-1)/\theta>0\) and \(L(t)=\mathbb E[e^{-tE}]\), the conditional
moments above give

\[
 L''(t)+\frac{2}{t+1/\theta}L'(t)-E_0^2 L(t)=0.                 \tag{3}
\]

Indeed the averaged second derivative has the two factors \(-cE/\theta\)
and \(c^2(E^2-E_0^2)/(2\theta^2)\). Set \(u=t+1/\theta\). Then
\((uL)''=E_0^2(uL)\). The bound \(0<L(t)\le e^{-tE_0}\) removes the
growing solution since \(J>0\). Normalization at zero gives

\[
 L(t)=\frac{e^{-tE_0}}{1+\theta t}.
\]

Uniqueness of the Laplace transform gives the stated law. At \(s=0\) we
have \(W=(c-1)E\), and the work identity gives (1). Differentiating (1)
in temperature gives
\((S(\theta,\omega)-S(\theta,c\omega))/k_B=\log c\). Taking \(c=2\)
contradicts condition 2. The familiar classical oscillator entropy
already has this dependence. The converse above is what prevents a
different positive distribution of Newtonian moving orbits from
repairing it while keeping the source work.
\href{https://arxiv.org/pdf/physics/0210128}{Boyer, oscillator
thermodynamics, §§III and VI} (passages).

\subsection{An action scale from two finite work
averages}\label{an-action-scale-from-two-finite-work-averages}

Set \(G(\theta,\omega)=F(\theta,2\omega)-F(\theta,\omega)\). Compare the
same compression at two positive temperatures using

\[
 B(\theta,\omega)=
 \frac{G(\theta,\omega)-G(\theta/2,\omega)}{\theta/2}.
                                                                    \tag{4}
\]

This is the average of the isothermal entropy loss in units of \(k_B\)
over the interval \([\theta/2,\theta]\) (or its value at an intermediate
temperature by the mean value theorem). It uses only two finite work
averages. An additive temperature-independent endpoint energy cancels.
For Newtonian motion (1) gives \(B=\log 2\). Nernst gives \(B\to0\), so
there is already a strict separation at a \emph{fixed positive}
temperature. Fix this temperature before refining the mechanical time
mesh or comparing trial actions.

To define an action directly from thermodynamics, keep the usual
mechanical source conjugacy
\(F_\omega=m\omega\langle r^2\rangle=U/\omega\) in equilibrium, with
\(U=F-\theta F_\theta\). The last mechanical equality is the harmonic
virial law at fixed \(J\). Under precisely this work law oscillator
entropy has Wien's scaling form, a function of \(\theta/\omega\) at
fixed body and sector. Thus (4) is a continuous function of
\(x=\theta/(2\omega)\), which has action units.
\href{https://arxiv.org/pdf/physics/0210128}{Boyer, equations (1), (8)
and (13)} (passages). Assume classical thermal correspondence
\(B(x)\to\log 2\) at large \(x\). For existence of the scale it suffices
that \(B\) exceed \(\log(3/2)\) at one finite \(x\). Then

\[
 A_*=\inf\{x>0:B(x)\ge\log(3/2)\},\qquad 0<A_*<\infty .          \tag{5}
\]

Indeed there is an entire neighborhood of zero below this threshold, and
the defining set is nonempty. Continuity gives a first crossing even if
the response elsewhere is nonmonotone. The choice \(3/2\) is the
midpoint of the two endpoint ratios \(e^B=1,2\). Other intermediate
thresholds define other finite positive scales. No action-valued cell
size or absolute entropy enters (4) or (5).

\subsubsection{The reduced action in the regular canonical
theory}\label{the-reduced-action-in-the-regular-canonical-theory}

For a spinless regular planar oscillator with trial reduced action
\(a>0\) and \(J=ka\), \(k\) a positive integer, the exact radial
energies are \(E_n=\omega[J+(2n+1)a]\), \(n\ge0\), each simple at fixed
\(k\). The Helmholtz difference follows from this equally spaced
spectrum. \href{https://doi.org/10.3390/sym16030297}{Mazhar et al.,
equation (56)} (passage) and
\href{rivero-1998-conjecture-central-forces.md\#5-hooke-and-kepler-radial-actions-and-maslov-phases}{radial
period and action}. Writing \(z=2a\omega/\theta=a/x\) gives

\[
 \begin{aligned}
 G_a(\theta,\omega)&=\omega(J+a)+\theta\log(1+e^{-z}),\\
 e^{B_a}&=\frac{(1+e^{-z})^2}{1+e^{-2z}}=1+\operatorname{sech}z .
 \end{aligned}                                                       \tag{6}
\]

Thus any observed \(B\le b<\log 2\), \(0<b\), implies

\[
 a\ \ge\ \frac{\theta}{2\omega}
       \operatorname{arcosh}\!\left(\frac1{e^b-1}\right)>0.
                                                                    \tag{7}
\]

Equality in the measured response determines one trial action through
(6), subject to \(J/a\) integral. In this theory
\(a=A_*\operatorname{arcosh}2\). The vanishing-temperature and
vanishing-action limits in (6) give different answers in opposite
orders. Each positive \(a=J/k\) separately obeys Nernst. What excludes
an entire sequence \(a\downarrow0\) for \emph{one} preparation is the
shared finite-temperature value (4).

For example take \(\theta=J\omega/\log 2\). At \(a=J\) and \(a=J/2\),
respectively, \(B=\log(25/17)\) and \(\log(9/5)\). Both trials use the
same angular momentum, two source strengths, two temperatures and work
rule. Even a single-temperature bounded reading separates them. For the
\(s=0,c=2\) quarter-period change, \(W=E\) also holds for the two
quantum energy readings. The harmonic evolution interchanges position
and momentum as in (2), giving \(U_p^\dagger H_{2\omega}U_p=2H_\omega\).
On the initial thermal state the mean of \(e^{-(W-E_0)/\theta}\) is
\(1/[2\cosh(a\omega/\theta)]\). The two means are \(2/5\) and
\(\sqrt2/3\), compared with \(1/2\) for the Newtonian work law. Rounding
\(W-E_0\) down in bins of width \(\theta/256\) and giving the reading
zero at energies above \(8\theta\) changes each mean by at most
\(e^{1/256}-1+e^{-8}<1/128\). All three intervals are disjoint.

An exactly smooth compression with the same endpoint relation is also
available. Take a smooth scale factor \(b(t)\) constant at 1 initially
and \(1/\sqrt2\) finally. Use the standard harmonic dilation with
\(\Omega^2=\omega^2/b^4-\ddot b/b\). A sufficiently long finite duration
keeps \(\Omega^2>0\). The construction acts identically on both
Cartesian coordinates and preserves \(J\). The endpoint energies scale
by two. The same profile works for both trial actions. Initial and final
energy readings can lie anywhere in fixed time windows on the constant
plateaus. \href{https://arxiv.org/pdf/0910.0709}{Chen et al., Ermakov
equation and endpoint conditions} (passage, reversed compression).

The dimensional inputs here are \(m\), the prepared angular momentum
\(J\), the source frequencies and their durations in physical time, and
the thermal energies \(\theta,\theta/2\). \(k_B\) converts temperature
to energy. One may equivalently fix independent length and time units
\(L,\tau\) and set \(J=mL^2/\tau,\ \omega=1/\tau\). The finite energy
bins and plateau windows can then be fixed in those same units before
choosing \(a\). Both \(B\) and the threshold in (5) are dimensionless.
With the mechanical reference \(S_0=J\), the dimensionless trial is
\(\epsilon=a/J\). The physical trial \(a\) tends to zero with
\(\epsilon\) at fixed \(J\). The time mesh is a separate regulator. In
the canonical identification \(a=\hbar\) the ordinary Planck action is
\(h_P=2\pi a\).

\textbf{A radius reading with a fixed ruler and clock.} Take those same
independent \(L,\tau\) and the same thermal preparation at
\(\theta=J\omega/\log2\). Sample its initial stationary radius at a
uniform time in \([\tau,2\tau]\). Use bin width \(\rho=L/1024\) and
report \(f(r)=\min\{64L,\rho\lfloor r/\rho\rfloor\}^2\). The equilibrium
virial relation gives \(\langle r^2\rangle=U/(m\omega^2)\). The two
geometric level populations give \(U/(J\omega)=8/3,5/2\) at \(a=J,J/2\),
respectively. Their energy second moments in units of \((J\omega)^2\)
are \(80/9,33/4\).

For each regular oscillator eigenfunction \(\psi\), integration by parts
gives \(2\operatorname{Re}\langle K\psi,V\psi\rangle
\ge -a^2\omega^2\|\psi\|^2\), since \(\Delta V=2m\omega^2\) in the plane
and the remaining integral is \((a^2/m)\int V|\nabla\psi|^2\). Thus
\(\langle V^2\rangle\le\langle H_\omega^2\rangle+a^2\omega^2\) also in
the thermal mixtures. Both have \(\langle r^4\rangle\le40L^4\) and
\(\langle r\rangle<2L\). The loss from rounding and saturation is at
most \(2\rho\langle r\rangle+\langle r^4\rangle/(64L)^2<L^2/64\). The
intervals for \(\langle f\rangle/L^2\) therefore lie within
\([8/3-1/64,8/3]\) and \([5/2-1/64,5/2]\), respectively. The action
changes the thermal populations and their position laws. The same finite
spatial bins and the same time window act on both.

\subsection{Why positive angular momentum
matters}\label{why-positive-angular-momentum-matters}

\textbf{A classical exception at zero angular momentum.} In the
invariant diameter of the same central force take the common reference
measure \(dq\,dp/(2\pi b_0)+\delta_{(0,0)}\), where \(b_0>0\) has action
units. Every change of spring strength preserves area and fixes the
origin. The Gibbs weights on this measure therefore obey the work
identity by change of variables, for the full mechanical work. Its
partition function is \(1+\theta/(b_0\omega)\). Its entropy tends to the
same value at both strengths as \(\theta\downarrow0\). Its
high-temperature loss tends to \(k_B\log2\). It has continuous Newtonian
trajectories and no quantum phase. The extra weight gives a thermal
scale \(b_0\). Here (3) admits a constant term in \(L\), the mass at the
fixed point. The strictly rotating sector removed exactly that
possibility. Also, an entropy difference is insensitive to
temperature-independent shifts of level energies. An absolute zero-point
energy requires further input.

\subsection{Two short force changes on a Kepler
orbit}\label{two-short-force-changes-on-a-kepler-orbit}

Take two Kepler strengths \(\mu_2>\mu_1>0\). With \(v=\dot r\), write
\(H_\mu=mv^2/2+U_\mu(r)\), where

\[
 V(r)=-m/r,\qquad U_\mu(r)=\frac{J^2}{2mr^2}+\mu V(r),\qquad
 R_\mu=\frac{J^2}{m^2\mu},\qquad \nu_\mu=\frac{m^3\mu^2}{J^3}.
\]

\(R_\mu,\nu_\mu\) are the circular radius and its radial frequency. The
unconfined Kepler canonical partition function diverges at large radius.
Use the same two fixed elastic walls \(0<r_-<r_+\) for all strengths in
an open interval containing \(\mu_1,\mu_2\), with every \(R_\mu\)
strictly inside. The walls perform zero work. The effective radial force
points into the annulus at each wall, so a stationary Newtonian law
spends zero probability at the endpoints, including turning points with
zero speed. Assume finite mean kinetic energy.
\href{self-sufficiency-contrast.md\#passivity-under-cyclic-changes-of-the-attraction}{Kepler
actions and circular limit} and \href{dirichlet-wall-impulse.md}{elastic
and Dirichlet walls}.

\textbf{Theorem.} For this source, condition 1 and stationarity
determine both the radial probability \(P_\mu\) and the conditional mean
square radial speed.

\[
 P_\mu(dr)=\frac{e^{-U_\mu(r)/\theta}}{I_\mu(\theta)}\,dr,\qquad
 \mathbb E_\mu[v^2\mid r]=\theta/m,\qquad
 I_\mu(\theta)=\int_{r_-}^{r_+} e^{-U_\mu(r)/\theta}\,dr .           \tag{8}
\]

Only the scalar source work and the motion enter the proof. In
particular it also applies after marginalizing any hidden state that
leaves this source work and equation of motion intact. The established
work-weighted endpoint formula starts with a Gibbs ensemble. Here two
strength changes and their short-time motion determine the radius
marginal of an initially unspecified stationary law.
\href{https://doi.org/10.1073/pnas.071034098}{Hummer and Szabo, equation
(4)} (passage, forward identity).

\textbf{Proof.} Change \(\mu\) to \(\lambda=\mu+\delta\), let the orbit
move for a time \(t\), then change to \(\alpha\). Its work is
\(\delta V(r_0)+(\alpha-\lambda)V(r_t)\). Comparing this route with
\(t=0\) in condition 1 shows that under the tilted initial measure
\(h_\delta(r_0)d\rho_\mu\), where \(h_\delta=e^{-\delta V/\theta}\) and
\(\rho_\mu\) is the whole stationary law, \(V(r_t)\) and \(V(r_0)\) have
identical Laplace transforms as \(\alpha\) ranges near \(\lambda\).
These variables are bounded. Equality on an interval determines their
laws, and \(V\) is one-to-one. Thus the \emph{radius marginal} of the
tilted measure is stationary under evolution at \(\lambda\).

The switches can be obtained with smooth ramps. The work throughout is
bounded by the total variation of the strength times \(\sup|V|\). For an
initial point in the interior the intermediate energy equals the
potential at either wall for at most two values of \(\lambda\), since
\(V(r_0)\ne V(r_\pm)\). By Fubini, for almost every intermediate
strength these grazing energies have zero initial probability. For those
strengths the limiting position flow is continuous through transverse
wall impacts. Bounded convergence therefore gives the stated work
comparison. Two distinct nonzero shifts from this full-measure set
suffice below.

Let \(\Pi(dr)=\mathbb E_\mu[v^2{\bf1}_{\{r\in dr\}}]\), a finite
measure. For a smooth test \(f\) supported in the interior, stationarity
of \(vf(r)\) at the original strength gives
\(\int f'\,d\Pi=m^{-1}\int f U_\mu'\,dP_\mu\). For the tilted measure
differentiate its constant mean of \(f(r_t)\) twice at zero, using
Newton's short-time acceleration at strength \(\lambda\). This gives

\[
 0=\int h_\delta f''\,d\Pi
    -\frac1m\int h_\delta f'(U_\mu'+\delta V')\,dP_\mu.             \tag{9}
\]

The compact support removes the wall impulse from both identities. The
second derivative along the flow is bounded by a constant times
\(1+v_0^2\), since the potentials are bounded on the annulus. Subtract
the original identity tested with \(h_\delta f'\). It follows that
\(h_\delta V'(\Pi-\theta P_\mu/m)\) has zero distributional derivative
and hence is constant times Lebesgue measure on the connected interior.
Two distinct shifts have a nonconstant ratio of tilts. Both constants
must vanish. Since \(V'>0\), \(\Pi=\theta P_\mu/m\). The original
identity now reads \(P_\mu'=-(U_\mu'/\theta)P_\mu\). Multiplication by
\(e^{U_\mu/\theta}\) and normalization give (8).

An immediate strength change then gives
\(G_K(\theta)=F(\theta,\mu_2)-F(\theta,\mu_1)
=-\theta\log(I_{\mu_2}/I_{\mu_1})\). Near each circle \(U_{\mu_i}\) has
a unique nondegenerate minimum \(E_{c,i}\) of curvature \(m\nu_i^2\),
where \(\nu_i=\nu_{\mu_i}\). The leading Laplace formula gives
\(I_{\mu_i}\sim e^{-E_{c,i}/\theta}\sqrt{2\pi\theta/m}/\nu_i\). In
consequence,

\[
 G_K(\theta)=E_{c,2}-E_{c,1}+\theta\Lambda+o(\theta),\qquad
 B_K(\theta):=\frac{G_K(\theta)-G_K(\theta/2)}{\theta/2}
       \longrightarrow \Lambda=\log(\nu_2/\nu_1)>0 .              \tag{10}
\]

The ordinary entropy loss has the same limit in units of \(k_B\) because
the weighted Laplace asymptotic gives
\(\langle U_{\mu_i}-E_{c,i}\rangle/\theta\to1/2\) at both minima. No
wall length or phase-volume normalization enters \(\Lambda\). Condition
2 for these two strengths requires \(B_K\to0\). The conflict comes from
the local acceleration in (9) and the source work.
\href{https://dlmf.nist.gov/2.3.iii}{DLMF §2.3(iii), applied on each
side of the minimum} (passage).

\subsubsection{A finite Kepler lower
bound}\label{a-finite-kepler-lower-bound}

For explicit dimensional choices take the earlier \(m,L,\tau\) and
\(J=mL^2/\tau\). Use \(\mu_1=L^3/\tau^2,\ \mu_2=\sqrt2\,\mu_1\),
\(r_-=L/2\) and \(r_+=2L\). Then \(\nu_1=1/\tau\), \(\nu_2=2/\tau\) and
\(\Lambda=\log2\). Fix a positive \(\theta\) low enough that the
Newtonian value in (10) exceeds \(3\Lambda/4\) and the physical
preparation obeying Nernst has \(|B_K|<\Lambda/4\). Existence follows
from the two limits. A regular canonical realization at each fixed
admissible \(a>0\) is consistent with Nernst. The Dirichlet radial
operator on this finite interval has a simple ground and positive
excitation gap in the fixed angular sector. The selected temperature
belongs to the preparation.

Here is a computable lower bound at that fixed \(\theta\). In flat
radial measure, at \(a=J/k\) the two Dirichlet Hamiltonians have kinetic
part \(-a^2\partial_r^2/(2m)\) and potentials \(U_{\mu_i}-a^2/(8mr^2)\).
Put \(M=\max_i\sup|U_{\mu_i}'|\) on the annulus,
\(\eta=\min\{\Lambda/12,1/2\}\). For the comparison take
\(N=\max\{1,\lceil 6M(r_+-r_-)/(\eta\theta)\rceil\}\) and
\(d=(r_+-r_-)/N\le\eta\theta/(6M)\). Every matching canonical trial must
satisfy

\[
 a>a_0,\qquad
 a_0=\min\left\{
 r_-\sqrt{\frac{4\eta m\theta}{3}},
 \ \frac{\eta d}{6}\sqrt{\frac{m\theta}{4\pi}}\right\}>0 .         \tag{11}
\]

To see this, denote the canonical partition function by \(Z_{a,i}(u)\)
at thermal energy \(u\). Write \(\ell_u=a\sqrt{2\pi/(m\,u)}\) and
\(Q=a^2/(8mr_-^2)\). Dirichlet--Neumann bracketing on cells of length
\(d\), with the extrema of the potential on each cell, gives

\[
 e^{-Md/u}(1-\ell_u/d)
 \ \le\ \frac{\ell_u Z_{a,i}(u)}{I_{\mu_i}(u)}
 \ \le\ e^{(Md+Q)/u}(1+\ell_u/d).
\]

Indeed the constant-potential cell kinetic sums are bounded below by
\(d/\ell_u-1\) for Dirichlet and above by \(d/\ell_u+1\) for Neumann, by
the integral bounds for the decreasing Gaussian. If \(a\le a_0\), at
both \(u=\theta,\theta/2\) the logarithm of the middle ratio is bounded
in magnitude by \((Md+Q)/u+2\ell_u/d\le\eta\). The two strengths give
error at most \(2u\eta\) in their free-energy difference. The two
temperatures give error at most \(6\eta\le\Lambda/2\) in \(B_K\). This
is too small to span the stipulated strict separation. The partition
here bounds a trace. Its cell length \(d\) can be chosen once and is
independent of the time subdivision in Newton's construction.
\href{https://doi.org/10.1017/9781009218436.007}{Frank, Laptev and
Weidl, Dirichlet--Neumann bracketing} (chapter summary) and
\href{self-sufficiency-contrast.md\#passivity-under-cyclic-changes-of-the-attraction}{regular
planar radial operator}.

\subsection{Consequence for STATE}\label{consequence-for-state}

The heat theorem and the finite-time work identity exclude the
zero-action Newtonian trajectory description for a thermal rotating body
in Proposition X's central force. The thermal action (5) is positive
under virial scaling and thermal correspondence. On a confined Kepler
orbit (9) gives the exclusion directly from the short-time acceleration.
Inequalities (7) and (11) bound the reduced action of a regular
canonical realization at fixed physical inputs. Two uncoupled bodies can
have different positive trial actions and satisfy all these tests
jointly. Their work averages multiply, free energies add and both
entropy differences vanish on cooling. A common phase scale for
interacting bodies and an exclusion from the bare trajectory limit
remain open.
\end{document}
